\section{Conclusion}
The creation of \textbf{AyurSpace} is a key step in the digital preservation and
democratization of Ayurvedic knowledge. By executing a hybrid neuro-symbolic architecture,
this framework addresses the major problems facing current plant identification
tools---absence of medicinal context and risk of generative hallucinations. Experimental
findings (96.4\% taxonomy accuracy, 99.1\% contextual reasoning compliance) confirm the
effectiveness of decoupling visual recognition from semantic analysis. The algorithmic
formalization of the \textit{Dosha} assessment transforms subjective clinical intuition
into replicable digital logic, enabling smart health decisions based on ancient wisdom
merged with robust modern mobile processing.

\section{Limitations}
\begin{enumerate}
    \item \textbf{Network Dependency}: Both AI APIs require an active internet connection.
    \item \textbf{API Cost}: Usage-based pricing models limit high-volume production.
    \item \textbf{Dataset Bias}: Plant.id may under-represent rare Indian medicinal herbs.
    \item \textbf{LLM Hallucinations}: Mitigated but not fully eliminated by the hybrid
    pipeline.
\end{enumerate}

\section{Future Scope}
\begin{enumerate}
    \item \textbf{Edge AI Optimization}: Quantized MobileNet models (int8) for complete
    offline plant identification on device, resolving connectivity issues in rural areas.
    \item \textbf{Augmented Reality (AR) Interfaces}: ARCore overlays to display medicinal
    properties, alerts, and plant structures directly on the camera viewfinder.
    \item \textbf{Telemedicine Integration}: ``Vaidya Connect'' component linking user
    Dosha baselines and scan histories with licensed Ayurvedic practitioners' dashboards.
\end{enumerate}
